\section{Allianz}

\subsection{Unternehmensstruktur und Klassifizierung als Großkonzern} % TODO vielleicht als \paragraph{Klassifizierung als Großkonzern}
Bei der Allianz Gruppe -- im weiteren Verlauf der Arbeit als Allianz bezeichnet~-- handelt es sich um ein multinationales, in ca.
70 Staaten vertretenes Versicherungsunternehmen, das weltweit mehr als 156.000 Angestellte beschäftigt und etwa 97~Mio. Kunden betreut
\cite{allianzFactSheet}. Diese ist in die übergeordnete Muttergesellschaft Allianz \glsdisp{se}{SE (lat. \textit{Societas Europaea}, dt.
Europäische Gesellschaft)} und die ihr untergeordneten Tochtergesellschaften wie die Allianz Versicherungs-\glsdisp{ag}{AG (Aktiengesellschaft)},
Allianz Technology \gls{se}, Allianz Global Investors \glsdisp{gmbh}{GmbH (Gesellschaft mit beschränkter Haftung)} etc. organisiert
\cite{allianzKonzern,allianz2025structure}. Durch diese Struktur gilt die Allianz laut § 18 des \glsdisp{aktg}{Aktiengesetzes (AktG)}
\cite{aktg} als Konzern und die Tochtergesellschaften entsprechend als Konzernunternehmen.

Ergänzend lässt sich die Allianz als Großkonzern klassifizieren. Nach § 267 \gls{hgb} \cite{hgb} ist ein
Großunternehmen durch die Erfüllung von zwei oder drei folgender Kriterien definiert:
\begin{enumerate}
    \item Der Jahresmittelwert der Mitarbeiteranzahl beträgt mindestens 251.
    \item Die Bilanzsumme übertrifft 25 Mio.~€.
    \item Die Umsatzerlöse der vergangenen zwölf Monate übertreffen 50 Mio.~€.
\end{enumerate}

Mit der signifikanten Anzahl an Beschäftigten sowie einem Umsatz von 186,9~Mrd.~€ \cite{geschäftsbericht2025,allianzFactSheet} und einer
Bilanzsumme von 1.024~Mrd.~€ \cite{geschäftsbericht2025}, die deutlich über die im \gls{hgb} festgelegten Grenzen hinausgehen, fällt
die Allianz somit unter die Kategorie der großen Kapitalgesellschaften.

\paragraph{Unternehmensorgane}
Darüber hinaus ist die interne Struktur der Allianz und ihrer Gesellschaften durch unterschiedliche Unternehmensorgane bestimmt, die abhängig
von der Rechtsform gesetzlich vorgeschrieben sind. Die Muttergesellschaft Allianz \gls{se} folgt dem in den §§~76--94, 95--116, 118--149
\gls{aktg} vorgegebenen Modell für Aktiengesellschaften \cite{aktg} und ist in die nachstehenden Unternehmensorgane gegliedert
\cite{allianzOrgane}: % TODO Prof fragen, ob \item auch zitiert sein soll mit \cite{allianzOrgane}
\begin{enumerate}
    \item Vorstand
    
    Für die Leitung der Gesellschaft ist ausschließlich der Vorstand zuständig, der gegenüber den übrigen zwei Organen berichtspflichtig ist.
    \item Aufsichtsrat
    
    Der Aufsichtsrat übernimmt die Aufgabe der Beaufsichtigung des Vorstands und kann diesen konsultativ unterstützen, beteiligt sich allerdings
    nicht unmittelbar an den Tätigkeiten des Vorstands. Zudem können durch den Aufsichtsrat die Vorstandsmitglieder bestellt und entlassen werden.
    \item Hauptversammlung
    
    Ebenfalls Bestandteil einer \gls{ag} ist die Hauptversammlung, zu der sämtliche Anteilseigner eingeladen werden, die Eigentümer mindestens
    einer Aktie des Unternehmens sind. Bei der Hauptversammlung werden unternehmensspezifische Entscheidungen getroffen.
\end{enumerate}

Während für die \glspl{ag} der Allianz das Vorhandensein der drei Unternehmensorgane im \gls{aktg} geregelt ist, sind Gesellschaften
anderer Rechtsform entsprechend der Gesetze strukturiert, die für diese Rechtsform gelten. So unterliegen die Tochtergesellschaften, die als eine
\gls{gmbh} organisiert sind, dem \gls{gmbhg} \cite{GmbHG}. Nach § 6 \gls{gmbhg} muss das Unternehmen durch mindestens einen Geschäftsführer
geleitet werden, der gemäß §§ 48--51 \gls{gmbhg} verpflichtet ist, Gesellschafterversammlungen zu veranstalten. Der Aufsichtsrat ist bei einer
\gls{gmbh} hingegen nicht obligatorisch \cite[§ 52]{GmbHG}.

Neben den aufgeführten Unternehmensorganen liegen in der Allianz \glspl{br} vor, die sich für die Bedürfnisse und Anliegen der Mitarbeiter einsetzen
und gemäß § 1 des \glsdisp{betrvg}{Betriebsverfassungsgesetzes (BetrVG)} \cite{BetrVG} durch die Angestellten gebildet werden können. Als Konzern
umfasst die Allianz einen \gls{kbr}, jeweils einen \gls{gbr} in der Muttergesellschaft und den Tochtergesellschaften sowie an einem Standort wie
München-Unterföhring einen örtlichen \gls{br} \cite{betriebsräteAllianz}.

Die Unternehmensorgane und \glspl{br} stehen in regelmäßigem Austausch miteinander, der teilweise staatlich sowie firmenintern vorgeschrieben ist
(siehe Abschnitt \ref{subsec:rechtliche_richtlinien}). Zur Kommunikation verwenden sie Werkzeuge, die den Angestellten der Allianz zur Verfügung
stehen. Diese werden im folgenden Abschnitt \ref{subsec:kommunikationsinfrastruktur} behandelt.

\paragraph{Hierarchieebenen} % TODO Prof fragen, wie Ebenen zitieren
Eine fortlaufende Kommunikation findet ebenso zwischen den Führungskräften verschiedener Ebenen und deren unterstellten Mitarbeitern statt. In der
Allianz existiert die hierarchische Struktur bestehend aus Referaten als unterste \gls{oe}, gefolgt von Abteilungen und Fachbereichen, über denen der
Vorstand bzw. die Geschäftsführung steht. Den Leitern dieser \glspl{oe}, dem Vorstand sowie den Mitarbeitern ohne Führungsfunktion ist jeweils
eine Ebene der sogenannten E-Hierarchie zugeordnet. Die Einteilung der Ebenen zu den entsprechenden Führungskräften ist in der vorliegenden Tabelle
\ref{tab:allianz_hierarchie} dargestellt, wobei es sich bei E0 um die höchste und bei E4 um die niedrigste Hierarchieebene handelt.

Die Tabelle verdeutlicht ebenfalls, welche agile Stellenarten (z.~B. Area Lead, Chapter Lead, Tribe Manager) einer klassischen Führungskraft in der
E-Hierarchie gleichgestellt sind. Da einige Inhaber einer agilen Rolle -- wie der Agile Coach -- keine Führungsposition ausüben \cite{spotify}, sich aber
dennoch an der Führungskräfte-Kommunikation beteiligen sollen, ist eine solche Zuordnung erforderlich. Allerdings können sowohl diese als auch die
hierarchische Struktur abhängig vom Konzernunternehmen voneinander abweichen, weshalb die Tabelle \ref{tab:allianz_hierarchie} lediglich eine
allgemeine Einteilung zeigt.

\hlinestable{tab:allianz_hierarchie}
{TODO Vergleich}
{l Q[wd=5.5cm,c] Q[wd=5.5cm,c]}
{
    Ebene & Klassische Führungskraft & Agile Stellenart \\
    E0 & \SetCell[c=2]{c} Vorstand / Geschäftsführer \\
    E1 & Fachbereichsleiter & Area Lead \\
    E2 & Abteilungsleiter & Senior Chapter Lead, Tribe Lead, Agile Coach u. a. \\
    E3 & Referatsleiter & Chapter Lead, Tribe Manager, Agile Master u. a. \\
    E4 & \SetCell[c=2]{c} Mitarbeiter ohne Führungsfunktion
}
[H]

Hinsichtlich der E-Hierarchie erfolgt zusätzlich eine Differenzierung zwischen den Führungskräften der ersten und zweiten Ebene, die als leitende
Führungskräfte gelten, und den nicht-leitenden E3-Angestellten.

\subsection{Kommunikationsinfrastruktur}
\label{subsec:kommunikationsinfrastruktur}
Für die Kommunikation im Arbeitsalltag, die Führungskräfte-Kommunikation und weitere Kommunikationsarten, die im Kapitel \ref{sec:theoretischer_hintergrund}
genauer behandelt werden, verwenden die Mitarbeiter der Allianz-Gesellschaften konzernweit standardisierte Werkzeuge. Diese werden im folgenden
Abschnitt vorgestellt.

\subsubsection{Microsoft-Werkzeuge und Cisco WebEx} % TODO ändern
Die Kommunikationslandschaft ist primär durch die von \gls{ms} bereitgestellten Kommunikationstools bestimmt. Mit der Konzernbetriebsvereinbarung aus
2019 \cite{KBVoffice365} wurde \gls{ms} Office 365 als Standard in der Allianz etabliert, sodass jeder Beschäftigte, der mit einem entsprechenden Konto
ausgestattet ist, \gls{ms} Teams für den regelmäßigen Nachrichtenaustausch und Outlook als E-Mail-Dienst einsetzen kann.

In Outlook sind E-Mail-Verteiler (im Folgenden: Verteiler) vorhanden, die der Gruppenkommunikation dienen. Hierbei werden durch die Mitarbeiter
überwiegend vorgefertigte Strukturverteiler \cite{strukturverteiler} -- mithilfe derer alle Personen einer \gls{oe} adressiert werden können
-- und Standortverteiler verwendet \cite{standortverteiler}. Diese beziehen die Angestellten ein, die in der jeweiligen Betriebsstätte arbeiten.

Strukturverteiler werden bspw. häufig zur Organisation von Mitarbeiterversammlungen innerhalb einer Abteilung oder eines Fachbereichs herangezogen, die
ggf. virtuell abgehalten werden. In diesem Fall kommt \gls{ms} Teams oder alternativ Cisco WebEx als weiteres Kommunikationswerkzeug zum Einsatz.

\subsubsection{Mailbroadcast} % TODO ändern
Über die bestehenden Outlook-Verteiler können jedoch nicht sämtliche gewünschte Empfängerkreise erreicht werden. Wenn exemplarisch ausgehend von einem
Standortverteiler ausschließlich die Inhaber einer agilen Rolle, die den leitenden Führungskräften gleichgestellt sind, benachrichtigt werden sollen,
ist das Herausfiltern korrekter Empfänger in Outlook nicht realisierbar. Daher haben die Mitarbeiter der Allianz, die eine solche Funktionalität benötigen,
die Möglichkeit, für die Anwendung Mailbroadcast berechtigt zu werden \cite{mbc2021,mbc2020}.

In Mailbroadcast liegen durch die Administratoren vorbereitete Verteiler vor \cite{mbc2020}, die u.~a. anhand der Struktureinheit und der E-Hierarchie
gefiltert werden können. Diese sind im Gegensatz zu Outlook-Verteilern für die Benutzer personalisiert, da sie basierend auf ihren Anliegen erstellt und
gepflegt werden. Des Weiteren kann innerhalb von Mailbroadcast eine Mail an die ermittelten Empfänger versendet werden oder das Ergebnis des Filtervorgangs
zwecks einer weiteren Verarbeitung -- z.~B. in Outlook -- als CSV-Datei exportiert werden \cite{mbc2021}.

\subsubsection{EverBridge}
Während die Microsoft-Kommunikationswerkzeuge und Mailbroadcast überwiegend in der regelmäßigen Unternehmenskommunikation\footnote{Die theoretischen
Grundlagen der regelmäßigen Unternehmenskommunikation sowie der Krisen- und Notfallkommunikation werden im Kapitel \ref{sec:theoretischer_hintergrund}
behandelt.} eingesetzt werden, wird bezüglich der anlassbezogenen Krisen- und Notfallkommunikation in der Allianz das Kri\-sen\-ma\-nagement-Tool EverBridge360
(im Folgenden: EverBridge) verwendet \cite{everbridgeInAllianz}. Über dieses System werden im unwahrscheinlichen Fall einer Krisensituation wie etwa eines
Erdbebens alle betroffenen Mitarbeiter in Form einer SMS-Nachricht, einer E-Mail oder eines Anrufs zeitnah benachrichtigt \cite{everbridgeInAllianz}.

Somit dient EverBridge der initialen Warnung in Risikofällen. Obgleich der Einsatz weiterer Systeme in der Notfallkommunikation derzeit nicht
geplant ist, da EverBridge konzernweit zur Anwendung kommt \cite{everbridgeInAllianz}, eignet sich Mailbroadcast für die nachfolgende Kommunikation,
in der die detaillierten Informationen zur jeweiligen Krisensituation per E-Mail über Mailbroadcast versendet werden.

Daher muss sichergestellt werden, dass Mailbroadcast den Anforderungen einer zuverlässigen Notfallkommunikation entspricht. Die Definition dieser
Anforderungen, deren Erfüllung durch den Ist-Zustand von Mailbroadcast sowie die daraus ggf. resultierende erforderliche Weiterentwicklung der
Anwendung werden im Abschnitt \ref{subsec:mbc_notfall} analysiert. % TODO \ref{} überprüfen im Laufe / am Ende der Bachelorarbeit

\subsection{Zentrale Personalsysteme der Allianz-Gesellschaften}
Die technische Grundlage für 

\subsubsection{Allianz Personalsystem (APAS-SAP)}
\subsubsection{Global Identity Access Management (GIAM)}
\subsubsection{SAP SuccessFactors}